\appendix
\section{ELBO}
\label{appendix:elbo_deriviation}

\newcommand{\Expec}{\mathbb{E}_{q}}

Our goal is to maximize the likelihood $p(\boldsymbol{y} \mid \boldsymbol{x})$. As SSL models use layer-aggregation methods for hidden representations from multiple layers, we need to derive a lower bound for the prior downstream task layer distribution.

For simplicity, let us denote the approximated posterior as $q_\theta(l\mid \boldsymbol{x})$, the true posterior which is a particular downstream task distribution as $p(l\mid \boldsymbol{x}, y)$, the prior layer distribution as  $p(l)$ and $\mathbb{E}_{l \sim q_\theta(l \mid \boldsymbol{x})}$ as $\Expec$~. Our goal is to approximate the posterior distribution so that it is maximally close to the true posterior. For that, let us derive a formula for KL-divergence:

\begin{equation*}
\begin{aligned}
    & \operatorname{KL}\left[ q_\theta(l \mid \boldsymbol{x}) \, \Vert \, p(l \mid \boldsymbol{x}, \boldsymbol{y})\right] \\
    &= \mathbb{E} \log \frac{q_\theta(l \mid \boldsymbol{x})}{p(l \mid \boldsymbol{x}, \boldsymbol{y})} \\
    &= \mathbb{E} \log \frac{q_\theta(l \mid \boldsymbol{x}) \cdot p(\boldsymbol{x}, \boldsymbol{y})}{p(l, \boldsymbol{y}, \boldsymbol{x})} \\
    &= \mathbb{E} \log \frac{q_\theta(l \mid \boldsymbol{x}) \cdot p(\boldsymbol{y} \mid \boldsymbol{x}) \cdot p(\boldsymbol{x})}{p(l, \boldsymbol{y} \mid \boldsymbol{x}) \cdot p(\boldsymbol{x})} \\
    &= \mathbb{E} \log \frac{q_\theta(l \mid \boldsymbol{x}) \cdot p(\boldsymbol{y} \mid \boldsymbol{x})}{p(l, \boldsymbol{y} \mid \boldsymbol{x})} \\
    &= \mathbb{E} \log \frac{q_\theta(l \mid \boldsymbol{x})}{p(l, \boldsymbol{y} \mid \boldsymbol{x})} + \mathbb{E} \log p(\boldsymbol{y} \mid \boldsymbol{x}) \\
    &= \mathbb{E} \log \frac{q_\theta(l \mid \boldsymbol{x})}{p(l, \boldsymbol{y} \mid \boldsymbol{x})} + \log p(\boldsymbol{y} \mid \boldsymbol{x}).
\end{aligned}
\end{equation*}


At this point, we still don't know how to access $p(l, \boldsymbol{y} \mid \boldsymbol{x})$. Let us rearrange the terms and write the log-likelihood $\log p(\boldsymbol{y} | \boldsymbol{x})$ as:

\begin{equation*}
\begin{aligned}
    \log p(\boldsymbol{y} \mid \boldsymbol{x}) &= \mathbb{E} \log \frac{p(l, \boldsymbol{y} \mid \boldsymbol{x})}{q_\theta(l \mid \boldsymbol{x})} + \operatorname{KL} \left(q_\theta(l \mid \boldsymbol{x}) \, \Vert \, p(l \mid \boldsymbol{x}, \boldsymbol{y})\right).
\end{aligned}
\end{equation*}


As log-likelihood is the logarithm of the probability, it is $\geq 0$. KL-divergence is a measure of distance which is also $\geq 0$, then the lower bound for likelihood can be written as:

\begin{equation*}
    \log p(\boldsymbol{y} \mid \boldsymbol{x}) \geq LB= \mathbb{E} \log \frac{p(l, \boldsymbol{y} \mid \boldsymbol{x})}{q_\theta(l
    \mid \boldsymbol{x})}
\end{equation*}

Our initial goal was to find a training objective for likelihood maximization. The training objective is the negative log-likelihood $-L(\theta) = LB$. Let us further simplify LB to bring it to a form which can be used to update gradients:

\begin{equation*}
\begin{aligned}
    LB &= -L(\theta) = \mathbb{E} \log \frac{p(l, \boldsymbol{y} \mid \boldsymbol{x})}{q_\theta(l \mid \boldsymbol{x})} = \mathbb{E} \log \frac{p(\boldsymbol{y} \mid \boldsymbol{x}, l) \cdot p(l)}{q_\theta(l \mid\boldsymbol{x})} \\
    &= \mathbb{E} \log p(\boldsymbol{y} \mid \boldsymbol{x}, l) + \mathbb{E} \log \frac{p(l)}{q_\theta(l \mid\boldsymbol{x})}.
\end{aligned}
\label{eq:lower_bound}
\end{equation*}

 % First, we need to express our LB through KL-divergence
From the definition of KL-divergence, $\operatorname{KL}\left[q_\theta(l \mid \boldsymbol{x}) \, \Vert \, p(l)\right]=\mathbb{E} \log \frac{q_\theta(l \mid \boldsymbol{x})}{p(l)}$, we get that\\
$\mathbb{E} \log \frac{p(l)}{q_\theta(l \mid \boldsymbol{x})}=-\operatorname{KL}\left(q_\theta(l \mid \boldsymbol{x}) \, \Vert \, p(l)\right)$. Which gives us equation~\ref{eq:elbo-pre-final-form} that is the same as equation~\ref{eq:objective}:

\begin{equation}
\label{eq:elbo-pre-final-form}
    L(\theta)=-\mathbb{E} \log p(\boldsymbol{y} \mid \boldsymbol{x}, l) + \operatorname{KL}\left(q_\theta(l \mid \boldsymbol{x}) \, \Vert \, p(l)\right).
\end{equation}

By expanding $\mathbb{E}$ and selecting the downstream model, $p_\theta$, we get equation~\ref{eq:elbo-final-form}:

\begin{equation}
\begin{aligned}
\label{eq:elbo-final-form}
    L(\theta) &= -\sum_{i=1}^n q_\theta(l=i \mid \boldsymbol{x}) \cdot \log p_\theta(\boldsymbol{y} \mid l=i, \boldsymbol{x}) \\ 
    &+ \operatorname{KL}\left(q_\theta(l \mid \boldsymbol{x}) \, \Vert \, p(l)\right).
\end{aligned}
\end{equation}
